% ===== Main text: append to the '(iii) Modifiers' paragraph in Data Curation =====
Before use, all $80$ modifiers were audited by one of the authors for semantic preservation and axis fidelity: $80/80$ preserved the original decision problem, and $76$--$79$ of $80$ ($95.0$--$98.8\%$, depending on whether three borderline cases are counted as errors) instantiate their labelled axis. The full protocol, the flagged items, and their bearing on the results are given in App.~\ref{app:modval}.

% ===== Limitations entry (before 'Sample diversity') =====
\textbf{Modifier validation is single-rater.} The audit of all $80$ modifiers (App.~\ref{app:modval}) was performed by one of the authors, not by independent annotators, so we cannot report an inter-annotator agreement statistic and the judgement of the borderline cases is a single person's. A multi-annotator validation would give a stronger guarantee that each modifier instantiates its intended axis.\\[2pt]

% ===== New appendix section (place before Schwartz Value Profile Enumeration) =====

\section{Modifier Validation Audit}
\label{app:modval}

\paragraph{Protocol.} All $80$ modifiers ($10$ scenarios $\times$ $8$ axes) were audited by one of the authors against two criteria, applied to every item rather than to a sample:

\begin{enumerate}[leftmargin=*, noitemsep, topsep=2pt]
  \item \textbf{Semantic preservation.} Does the sentence fit its scenario, add genuine contextual pressure, leave both actions viable, and avoid introducing a covert third option or restating an action label verbatim?
  \item \textbf{Axis fidelity.} Does the text actually instantiate the modifier axis it is labelled with?
\end{enumerate}

\noindent This is a single-rater audit by an author, not an independent multi-annotator study; no inter-annotator agreement statistic is therefore available, and we report it as an internal quality check on dataset construction rather than as external validation (see Limitations).

\paragraph{Semantic preservation: $80/80$ ($100\%$).} Every modifier is a coherent contextual addition to its scenario. None introduces a third option and none repeats an action label verbatim. Two items are somewhat redundant with their scenario premise --- \texttt{MOD\_SC003\_1\_02} and \texttt{MOD\_SC003\_2\_02} (\texttt{resource\_scarcity}) essentially restate the ``only one award'' condition already present in SC003\_1/SC003\_2 --- but they remain semantically valid and were retained.

\paragraph{Axis fidelity: $76$--$79$ of $80$ ($95.0$--$98.8\%$).} One item is a clear misfit and three are borderline, so we report both a strict and a lenient count rather than a single contested number. Table~\ref{tab:modval} lists every flagged item; all four were retained unchanged, so the reported flip rates include them.

\begin{table}[h]
\centering\footnotesize
\setlength{\tabcolsep}{3.5pt}
\begin{tabular}{@{}llp{3.1cm}@{}}
\toprule
\textbf{Item} & \textbf{Labelled} & \textbf{Issue} \\
\midrule
\texttt{SC003\_2\_01} & self\_pres. & Clear misfit: reads as benevolence/fairness, not self-interest \\
\addlinespace
\texttt{SC002\_1\_01} & self\_pres. & Borderline: ``mistakes more visible'' leans social\_visibility \\
\addlinespace
\texttt{SC002\_1\_04} & in\_out\_group & Borderline: overlaps the adoption dilemma rather than pure identity \\
\texttt{SC002\_2\_04} & in\_out\_group & Borderline: as above \\
\bottomrule
\end{tabular}
\caption{The four modifiers flagged for axis fidelity. Counting all four as errors gives $76/80$ ($95.0\%$); counting only the clear misfit gives $79/80$ ($98.8\%$).}
\label{tab:modval}
\end{table}

The remaining five axes --- \texttt{social\_visibility}, \texttt{time\_pressure}, \texttt{diffused\_responsibility}, \texttt{competence\_uncertainty}, and \texttt{authority\_signal} --- are $10/10$ each, and \texttt{resource\_scarcity} is clean on axis fidelity.

\paragraph{Bearing on the results.} The flagged items do not drive the paper's findings. The clear misfit and one borderline case sit on \texttt{self\_preservation}, one of the two dominant axes; both belong to scenarios (SC003\_2, SC002\_1) that are \emph{not} the concentration cell reported in \S\ref{sec:analysis}, and \texttt{self\_preservation} remains the largest positive McNemar effect in every model with them included. The two \texttt{in\_out\_group} items belong to an axis that is consistently mid-ranked and never significant in the direction of a pressure effect, so their classification does not affect any reported comparison.
